% Part: applied-modal-logics
% Chapter: epistemic-logic
% Section: bisimulations

\documentclass[../../../include/open-logic-section]{subfiles}

\begin{document}

\olfileid[hi]{aml}{el}{bsd}

\olsection{द्विसिमुलेशन}

हमारी ज्ञानात्मक भाषा की अभिव्यक्ति-शक्ति के बारे में एक शेष प्रश्न
प्रतिरूपों और उनमें सत्य सूत्रों के बीच के संबंध से जुड़ा है। फ़्रेम-अनुरूपता
के परिणामों से हमने देखा है कि जब कुछ सूत्र किसी फ़्रेम में मान्य होते हैं,
तो वे यह भी सुनिश्चित करते हैं कि वह फ़्रेम कुछ निश्चित गुणों को संतुष्ट
करे। किंतु क्या हमारी मोडल भाषा, उदाहरणार्थ, ऐसे जगत जिसमें एक स्वपरावर्ती
तीर है और जगतों की ऐसी अनंत शृंखला जिसमें प्रत्येक जगत अगले तक जाता है, के
बीच अंतर कर सकती है? अर्थात क्या कोई ऐसा सूत्र~$A$ है जो इन दोनों में से
केवल एक जगत पर सत्य हो?

द्विसिमुलेशन संबंधात्मक प्रतिरूपों के बीच परिभाषित किया जाने वाला ऐसा संबंध
है जो बताता है कि उनकी संरचना प्रभावतः एक ही है। जैसा हम देखेंगे, यह
प्रतिरूपों के बीच समतुल्यता की उस धारणा को ग्रहण करता है जिसे हमारी
ज्ञानात्मक भाषा में व्यक्त किया जा सकता है।

\begin{defn}[द्विसिमुलेशन]
मान लें कि $M_1 = \tuple{W_1, R_1, V_1}$ और
$M_2 = \tuple{W_2, R_2, V_2}$ दो संबंधात्मक प्रतिरूप हैं, और
$\mathcal{R} \subseteq W_1 \times W_2$ एक द्विपदी संबंध है। हम
$\mathcal{R}$ को \emph{द्विसिमुलेशन} कहते हैं, यदि प्रत्येक
$\tuple{w_1, w_2} \in \mathcal{R}$ के लिए निम्न शर्तें पूरी हों:

\begin{enumerate}
  \item प्रत्येक !!{propositional variable}~$p$ के लिए $w_1 \in V_1(p)$ तभी
  जब $w_2 \in V_2(p)$।
  \item सभी कर्ताओं $a \in A$ और जगतों $v_1 \in W_1$ के लिए, यदि
  $R_{1_a} w_1 v_1$, तो कोई $v_2 \in W_2$ ऐसा है कि
  $R_{2_a} w_2 v_2$ और $\tuple{v_1, v_2} \in \mathcal{R}$।
  \item सभी कर्ताओं $a \in A$ और जगतों $v_2 \in W_2$ के लिए, यदि
  $R_{2_a} w_2 v_2$, तो कोई $v_1 \in W_1$ ऐसा है कि
  $R_{1_a} w_1 v_1$ और $\tuple{v_1, v_2} \in \mathcal{R}$।
\end{enumerate}

जब $M_1$ और $M_2$ के बीच कोई द्विसिमुलेशन जगतों $w_1$ और $w_2$ को जोड़ता
है, तब हम $\tuple{M_1, w_1} \leftrightarroweq \tuple{M_2, w_2}$ भी लिखते
हैं और $\tuple{M_1, w_1}$ तथा $\tuple{M_2, w_2}$ को
\emph{द्विसिमुलेशन-समान} कहते हैं।
\end{defn}

द्विसिमुलेशन संबंध के अलग-अलग उपवाक्य अलग बातें सुनिश्चित करते हैं। उपवाक्य
1 यह सुनिश्चित करता है कि द्विसिमुलेशन-समान जगत समान मोडल-मुक्त सूत्रों को
संतुष्ट करेंगे, क्योंकि वह सभी !!{propositional variable}ों पर सहमति सुनिश्चित
करता है। अन्य दो उपवाक्य, जिन्हें क्रमशः ``आगे'' और ``पीछे'' कहते हैं,
सुनिश्चित करते हैं कि अभिगम्यता संबंधों की संरचना समान होगी।

\begin{thm}
यदि $\tuple{M_1, w_1} \leftrightarroweq \tuple{M_2, w_2}$, तो प्रत्येक
!!{formula}~$!A$ के लिए $\mSat{M_1}{!A}[w_1]$ तभी जब
$\mSat{M_2}{!A}[w_2]$।
\end{thm}

\begin{figure}
  \begin{center}
    \begin{tikzpicture}[modal]
      \node[world] (w1) {$w_1$}; 
      \node[world] (w2) [above left=of w1]{$w_2$}; 
      \node[world] (w3) [above right=of w1] {$w_3$};
      \draw[<->] (w1) to node {$a$} (w2);
      \draw[->,loop below] (w1) to node {$a$} (w1);
      \draw[<->] (w1) to [swap] node {$a$} (w3);
      \draw[->,loop above] (w2) to node {$a$} (w2) ;
      \draw[->,loop above] (w3) to node {$a$} (w3) ;
      
      \node[world](v1)[right of=w1, xshift=1.5in]{$v_1$};
      \node[world](v2)[above of=v1]{$v_2$};
      \draw[<->] (v1) to node {$a$} (v2);
      \draw[->,loop below] (v1) to node {$a$} (v1);
      \draw[->,loop above] (v2) to node {$a$} (v2) ;

      \draw[-,dotted] (w1) to (v1) ;
      \draw[-,dotted,bend left=45] (w2) to (v2) ;
      \draw[-,dotted,bend left=30] (w3) to (v2) ;
    \end{tikzpicture}
  \end{center}
  \caption{दो द्विसिमुलेशन-समान प्रतिरूप।}
  \ollabel{fig:bisimilar}
\end{figure}

यद्यपि \olref{fig:bisimilar} में चित्रित दोनों प्रतिरूप एक-दूसरे के बिल्कुल
समान नहीं हैं, फिर भी जगतों $w_1$ और~$v_1$ को जोड़ने वाला द्विसिमुलेशन है।
यह द्विसिमुलेशन $w_2$ और $w_3$ दोनों को~$v_2$ से भी जोड़ेगा। आशय यह है कि
हमारी मोडल भाषा में ऐसा कुछ व्यक्त नहीं किया जा सकता जो वास्तव में उनके बीच
अंतर करे। किंतु यदि $w_2$ और~$w_3$ भिन्न !!{propositional variable}ों को संतुष्ट
करते, तो स्थिति अलग होती।

\end{document}
